% GENERATED FILE. Edit canonical lessons, then run npm run generate:books.
% canonical-source-hash: fnv1a64:308d1faa2bd73cac
% canonical-lessons: MR-C515-abhiyanta, MR-C515-sallagar, MR-C515-anuvadak, MR-C515-lekhapal, MR-C515-atopne

\chapter{Experts at Work}
\label{ch:mr-experts-at-work}
\hlchaptermodality{515}

\begin{chapteropening}
\textbf{By the end of this chapter:} \emph{I can name some experts and say how they help.}

Complete the last lesson of chapter 515: I can name some experts and say how they help.
\end{chapteropening}

\section[\texorpdfstring{\mr{अभियंता} (abhiyantā)}{अभियंता (abhiyantā)}]{\mr{अभियंता} (abhiyantā) — an engineer}

\label{lesson:MR-C515-abhiyanta}



\begin{quote}
Before the new one: say the Marathi for to get done, then the Marathi for to struggle.
\end{quote}

\subsection*{You'll want to know: \mr{अभियंता}}
\textbf{\mr{अभियंता}} — \emph{abhiyantā} — \textquotedblleft{}an engineer\textquotedblright{}.

\begin{quote}
\mr{माझी} \mr{मावशी} \mr{अभियंता} \mr{आहे} \mr{आणि} \mr{पूल} \mr{बांधते}.
\end{quote}

\begin{quote}
\emph{mājhī māvśī abhiyantā āhe āṇi pūl bāndhte.}
\end{quote}

\begin{quote}
My aunt is an engineer and builds bridges.
\end{quote}

Masculine.

One word for a man or a woman.

\begin{cousinweb}
\textbf{\mr{अभियंता}} (\emph{abhiyantā}) is a modern Marathi coinage built from Sanskrit parts: \emph{abhi-}, \textquotedblleft{}towards\textquotedblright{}, and \emph{yantṛ}, \textquotedblleft{}one who guides or restrains\textquotedblright{}, related to \emph{yantra}, \textquotedblleft{}a machine\textquotedblright{} (Monier-Williams).
\end{cousinweb}

\begin{grammarlens}[title={What you've built}]
One more word for describing what you lived and felt.
\end{grammarlens}

\subsection*{Your turn}
\begin{itemize}
\raggedright
  \item \emph{Say it:} \emph{abhiyantā}
  \item \emph{Say it:} \emph{abhiyantā}, once more
  \item \emph{Say it:} \emph{mājhī māvśī abhiyantā āhe āṇi pūl bāndhte.} — My aunt is an engineer and builds bridges.
  \item \emph{Recall:} say the Marathi for to get done, then the Marathi for to struggle, then say \emph{abhiyantā} again
\end{itemize}

\begin{tcolorbox}[breakable,colback=teal!4,colframe=teal!35!black,title={Before you move on}]
What does \mr{अभियंता} mean? (\textquotedblleft{}An engineer\textquotedblright{}.) Say it once more.
\end{tcolorbox}
\section[\texorpdfstring{\mr{सल्लागार} (sallāgār)}{सल्लागार (sallāgār)}]{\mr{सल्लागार} (sallāgār) — a consultant; an adviser}

\label{lesson:MR-C515-sallagar}



\begin{quote}
Before the new one: say the Marathi for to struggle, then the Marathi for an engineer.
\end{quote}

\subsection*{You'll want to know: \mr{सल्लागार}}
\textbf{\mr{सल्लागार}} — \emph{sallāgār} — \textquotedblleft{}a consultant; an adviser\textquotedblright{}.

\begin{quote}
\mr{कंपनीने} \mr{एका} \mr{सल्लागाराची} \mr{मदत} \mr{घेतली}.
\end{quote}

\begin{quote}
\emph{kampanīne ekā sallāgārācī madat ghetlī.}
\end{quote}

\begin{quote}
The company took the help of a consultant.
\end{quote}

Masculine.

\begin{cousinweb}
\textbf{\mr{सल्लागार}} (\emph{sallāgār}) is built on \mr{सल्ला}, \textquotedblleft{}advice\textquotedblright{}, a loan usually traced to Arabic \emph{ṣalāḥ}, with the Persian ending \emph{-gār}, \textquotedblleft{}one who does\textquotedblright{}.
\end{cousinweb}

\begin{grammarlens}[title={What you've built}]
One more word for describing what you lived and felt.
\end{grammarlens}

\subsection*{Your turn}
\begin{itemize}
\raggedright
  \item \emph{Say it:} \emph{sallāgār}
  \item \emph{Say it:} \emph{sallāgār}, once more
  \item \emph{Say it:} \emph{kampanīne ekā sallāgārācī madat ghetlī.} — The company took the help of a consultant.
  \item \emph{Recall:} say the Marathi for to struggle, then the Marathi for an engineer, then say \emph{sallāgār} again
\end{itemize}

\begin{tcolorbox}[breakable,colback=teal!4,colframe=teal!35!black,title={Before you move on}]
What does \mr{सल्लागार} mean? (\textquotedblleft{}A consultant; an adviser\textquotedblright{}.) Say it once more.
\end{tcolorbox}
\section[\texorpdfstring{\mr{अनुवादक} (anuvādak)}{अनुवादक (anuvādak)}]{\mr{अनुवादक} (anuvādak) — a translator}

\label{lesson:MR-C515-anuvadak}



\begin{quote}
Before the new one: say the Marathi for an engineer, then the Marathi for a consultant.
\end{quote}

\subsection*{You'll want to know: \mr{अनुवादक}}
\textbf{\mr{अनुवादक}} — \emph{anuvādak} — \textquotedblleft{}a translator\textquotedblright{}.

\begin{quote}
\mr{अनुवादकाने} \mr{इंग्रजी} \mr{पत्राचं} \mr{मराठीत} \mr{भाषांतर} \mr{केलं}.
\end{quote}

\begin{quote}
\emph{anuvādakāne iṅgrajī patrācaṁ marāṭhīt bhāṣāntar kelaṁ.}
\end{quote}

\begin{quote}
The translator turned the English letter into Marathi.
\end{quote}

Masculine.

\begin{cousinweb}
\textbf{\mr{अनुवादक}} (\emph{anuvādak}) is a tatsama, taken from Sanskrit \emph{anuvādaka}, \textquotedblleft{}repeating with explanation\textquotedblright{} (Monier-Williams).
\end{cousinweb}

\begin{grammarlens}[title={What you've built}]
One more word for describing what you lived and felt.
\end{grammarlens}

\subsection*{Your turn}
\begin{itemize}
\raggedright
  \item \emph{Say it:} \emph{anuvādak}
  \item \emph{Say it:} \emph{anuvādak}, once more
  \item \emph{Say it:} \emph{anuvādakāne iṅgrajī patrācaṁ marāṭhīt bhāṣāntar kelaṁ.} — The translator turned the English letter into Marathi.
  \item \emph{Recall:} say the Marathi for an engineer, then the Marathi for a consultant, then say \emph{anuvādak} again
\end{itemize}

\begin{tcolorbox}[breakable,colback=teal!4,colframe=teal!35!black,title={Before you move on}]
What does \mr{अनुवादक} mean? (\textquotedblleft{}A translator\textquotedblright{}.) Say it once more.
\end{tcolorbox}
\section[\texorpdfstring{\mr{लेखापाल} (lekhāpāl)}{लेखापाल (lekhāpāl)}]{\mr{लेखापाल} (lekhāpāl) — an accountant}

\label{lesson:MR-C515-lekhapal}



\begin{quote}
Before the new one: say the Marathi for a consultant, then the Marathi for a translator.
\end{quote}

\subsection*{You'll want to know: \mr{लेखापाल}}
\textbf{\mr{लेखापाल}} — \emph{lekhāpāl} — \textquotedblleft{}an accountant\textquotedblright{}.

\begin{quote}
\mr{लेखापाल} \mr{दर} \mr{महिन्याला} \mr{कंपनीचा} \mr{हिशेब} \mr{तपासतो}.
\end{quote}

\begin{quote}
\emph{lekhāpāl dar mahinyālā kampanīcā hiśeb tapāsto.}
\end{quote}

\begin{quote}
The accountant checks the company's accounts every month.
\end{quote}

Masculine.

\begin{cousinweb}
\textbf{\mr{लेखापाल}} (\emph{lekhāpāl}) is a modern tatsama coinage from Sanskrit \emph{lekhā}, \textquotedblleft{}a line, an account\textquotedblright{}, and \emph{pāla}, \textquotedblleft{}keeper\textquotedblright{} (Monier-Williams).
\end{cousinweb}

\begin{grammarlens}[title={What you've built}]
One more word for describing what you lived and felt.
\end{grammarlens}

\subsection*{Your turn}
\begin{itemize}
\raggedright
  \item \emph{Say it:} \emph{lekhāpāl}
  \item \emph{Say it:} \emph{lekhāpāl}, once more
  \item \emph{Say it:} \emph{lekhāpāl dar mahinyālā kampanīcā hiśeb tapāsto.} — The accountant checks the company's accounts every month.
  \item \emph{Recall:} say the Marathi for a consultant, then the Marathi for a translator, then say \emph{lekhāpāl} again
\end{itemize}

\begin{tcolorbox}[breakable,colback=teal!4,colframe=teal!35!black,title={Before you move on}]
What does \mr{लेखापाल} mean? (\textquotedblleft{}An accountant\textquotedblright{}.) Say it once more.
\end{tcolorbox}
\section[\texorpdfstring{\mr{आटोपणे} (āṭopṇe)}{आटोपणे (āṭopṇe)}]{\mr{आटोपणे} (āṭopṇe) — to wrap up; to finish}

\label{lesson:MR-C515-atopne}



\begin{quote}
Before the new one: say the Marathi for a translator, then the Marathi for an accountant.
\end{quote}

\subsection*{You'll want to know: \mr{आटोपणे}}
\textbf{\mr{आटोपणे}} — \emph{āṭopṇe} — \textquotedblleft{}to wrap up; to finish\textquotedblright{}.

\begin{quote}
\mr{लेखापालाने} \mr{संध्याकाळी} \mr{काम} \mr{आटोपलं} \mr{आणि} \mr{तो} \mr{घरी} \mr{गेला}.
\end{quote}

\begin{quote}
\emph{lekhāpālāne sandhyākāḷī kām āṭoplaṁ āṇi to gharī gelā.}
\end{quote}

\begin{quote}
The accountant wrapped up his work in the evening and went home.
\end{quote}

Past, with -\mr{ने} (he …): \mr{त्याने} \mr{काम} \mr{आटोपलं}, \emph{tyāne kām āṭoplaṁ}.

Close to \mr{उरकणे}: both say a job was finished.

\begin{cousinweb}
\textbf{\mr{आटोपणे}} (\emph{āṭopṇe}) is a Marathi verb for bringing a task to a close; its origin is not settled.
\end{cousinweb}

\begin{grammarlens}[title={What you've built}]
One more word for describing what you lived and felt.
\end{grammarlens}

\subsection*{Your turn}
\begin{itemize}
\raggedright
  \item \emph{Say it:} \emph{āṭopṇe}
  \item \emph{Say it:} \emph{āṭopṇe}, once more
  \item \emph{Say it:} \emph{lekhāpālāne sandhyākāḷī kām āṭoplaṁ āṇi to gharī gelā.} — The accountant wrapped up his work in the evening and went home.
  \item \emph{Recall:} say the Marathi for a translator, then the Marathi for an accountant, then say \emph{āṭopṇe} again
\end{itemize}

\begin{tcolorbox}[breakable,colback=teal!4,colframe=teal!35!black,title={Before you move on}]
What does \mr{आटोपणे} mean? (\textquotedblleft{}To wrap up; to finish\textquotedblright{}.) Say it once more.
\end{tcolorbox}
